Participants were stratified into five clinically defined diagnostic categories to reflect the full continuum of Alzheimer’s disease pathology: Cognitively Normal (CN; $n = 577$), Alzheimer’s disease dementia (AD; $n = 317$), and three subtypes of Mild Cognitive Impairment (MCI). To account for the heterogeneity of MCI, we further subclassified individuals based on delayed recall performance on the WMS-R Logical Memory II Story A task: Early MCI (EMCI; $n = 291$) denoted a milder impairment stage, whereas Late MCI (LMCI; $n = 152$) denoted a more severe impairment stage. The remaining MCI participants not classified as EMCI or LMCI were grouped as MCI-unclassified (MCI-U; $n = 520$). These categories define cross-sectional diagnostic and MCI-stage comparisons, without establishing longitudinal prognosis.
